\chapter{Carte des cultures}
\label{chap:carte}

\section{Vue d'ensemble}

L'écran \emph{Carte des cultures} (raccourci \texttt{Ctrl+9}) affiche l'occupation de chaque lieu sur les 12 mois de la saison. Chaque ligne représente un \textbf{lieu} (planche, parcelle, serre…) et chaque barre colorée représente une \textbf{plantation} positionnée selon ses dates :

\begin{itemize}
  \item Pour une plantation annuelle (\texttt{Cycle}), la barre s'étend de la première opération connue (semis, repiquage ou première récolte) jusqu'à la dernière récolte.
  \item Pour une plantation pluriannuelle (\texttt{Perennial}), la barre couvre toute l'année — la plantation occupe le lieu en permanence.
\end{itemize}

La couleur d'une barre est dérivée déterministiquement de la \textbf{famille botanique} de la culture, ce qui aide à repérer les rotations en un coup d'œil.

\section{Sélectionner une plantation}

Cliquer sur une barre la sélectionne (bordure renforcée et opacité maximale). Recliquer la désélectionne. Une seule plantation peut être sélectionnée à la fois. Une fois une plantation choisie, les boutons \textbf{Déplacer vers…}, \textbf{Diviser} et \textbf{Désélectionner} s'activent dans l'en-tête.

\section{Déplacer une plantation}

\begin{enumerate}
  \item Cliquez sur la plantation à déplacer.
  \item Cliquez sur \textbf{Déplacer vers…} dans l'en-tête.
  \item Une fenêtre apparaît avec la liste de tous les lieux. Cliquez sur celui de destination.
  \item La plantation passe immédiatement sur la nouvelle ligne. Ses tâches et ses récoltes annuelles restent attachées à elle (elles suivent la plantation, pas le lieu).
\end{enumerate}

\section{Diviser une plantation}

Une plantation peut être divisée entre deux lieux distincts — utile lorsque vous avez surévalué l'occupation d'une planche ou que vous voulez expérimenter une autre disposition.

\begin{enumerate}
  \item Cliquez sur la plantation à diviser.
  \item Cliquez sur \textbf{Diviser} dans l'en-tête.
  \item Une fenêtre apparaît avec deux blocs : \textbf{Part A} et \textbf{Part B}, pré-remplis avec un partage 50/50 (la part A garde le lieu courant, la part B propose le lieu suivant dans la liste).
  \item Ajustez les surfaces (m²) et le nombre de plants de chaque part.
  \item Validez avec \textbf{Diviser}.
\end{enumerate}

\subsection*{Ce qui se passe en interne}

\begin{itemize}
  \item La \textbf{Part A} \emph{conserve} la plantation d'origine (même identifiant). Toutes ses tâches et ses récoltes annuelles restent attachées à elle.
  \item La \textbf{Part B} \emph{est créée} comme une nouvelle plantation, qui hérite de la variété, des dates, du nom et des notes de l'originale. Elle démarre avec aucune tâche ni récolte associée.
\end{itemize}

\textbf{Note} : la somme des surfaces et le nombre total de plants ne sont pas vérifiés automatiquement par rapport à l'originale — Pomone vous laisse le contrôle. C'est à vous de choisir si vous voulez répartir strictement la surface initiale ou en profiter pour ajuster.

\begin{notebox}
\textbf{Prévu.} Déplacement d'une barre directement à la souris (drag-and-drop) ; division en trois parts ou plus en un seul écran ; alertes de rotation (interdire la même famille deux saisons consécutives sur la même planche).
\end{notebox}
